% PACLIC 40 uses ACL 2026 two-column style.
% Replace this preamble with the official ACL 2026 template files before submission.
\documentclass[11pt]{article}
\usepackage[margin=1in]{geometry}
\usepackage{times}
\usepackage{latexsym}
\usepackage{booktabs}
\usepackage{graphicx}
\usepackage{amsmath}
\usepackage{url}

\title{Online Prediction or Contextual Integration?\\Comparing Autoregressive Surprisal and Masked Pseudo-Surprisal in Human Reading}

\author{Anonymous PACLIC submission}

\begin{document}
\maketitle

\begin{abstract}
Language-model surprisal is widely used as a computational predictor of human reading difficulty, but different language-model objectives may correspond to different stages of processing. This paper compares local n-gram surprisal, left-to-right autoregressive transformer surprisal, and bidirectional masked pseudo-surprisal on word-level eye-tracking measures. We ask whether autoregressive surprisal better captures online prediction, while masked pseudo-surprisal captures later contextual compatibility. Using a naturalistic eye-tracking corpus, we evaluate early and late reading-time measures under lexical controls, incremental variance analysis, and sentence-level cross-validation. The study is designed as a stage-sensitive analysis rather than a model leaderboard.
\end{abstract}

\section{Introduction}
Write the motivation here. Main claim: the best computational predictor may depend on the eye-tracking measure.

\section{Related Work}
Cover surprisal theory, n-gram/local statistics, transformer surprisal and reading time, and pseudo-log-likelihood for masked models.

\section{Corpus and Preprocessing}
Describe the chosen corpus, participant subset, reading-time measures, filtering, token alignment, and exclusion rates.

\section{Predictors}
\subsection{Lexical controls}
Word length, log frequency, position in sentence, sentence length, sentence-initial and sentence-final indicators.

\subsection{N-gram surprisal}
Define the n-gram baseline.

\subsection{Autoregressive surprisal}
For a word $w_i$ in context $w_{<i}$:
\begin{equation}
S_{AR}(w_i) = -\log_2 P(w_i \mid w_{<i}).
\end{equation}

\subsection{Masked pseudo-surprisal}
For a masked model using left and right context:
\begin{equation}
S_{MLM}(w_i) = -\sum_{t \in T(w_i)} \log_2 P(x_t \mid x_{\setminus t}).
\end{equation}
Explain that this is contextual compatibility, not direct online prediction.

\section{Experiments}
List model ladder: lexical baseline, lexical+n-gram, lexical+AR, lexical+MLM, lexical+n-gram+AR, full model.

\section{Results}
Include a table like this:

\begin{table}[h]
\centering
\small
\begin{tabular}{lrrrr}
\toprule
Target & Model & $R^2$ & $\Delta R^2$ & CV RMSE \\
\midrule
Gaze duration & Lexical & -- & -- & -- \\
Gaze duration & + AR & -- & -- & -- \\
Total time & Lexical & -- & -- & -- \\
Total time & + MLM & -- & -- & -- \\
\bottomrule
\end{tabular}
\caption{Model comparison under lexical controls.}
\end{table}

\section{Discussion}
Interpret whether AR surprisal is stronger for earlier measures and whether masked pseudo-surprisal adds variance for later measures.

\section{Limitations and Ethics}
Mention corpus genre, tokenization mismatch, masked pseudo-surprisal interpretation, and use of public/de-identified eye-tracking data.

\section{Conclusion}
Summarise the stage-sensitive conclusion.

\bibliographystyle{acl_natbib}
\bibliography{references}

\end{document}
